\documentclass[../../main.tex]{subfiles}
\begin{document}

{\bf[解]}
与えられた条件
\begin{align}
  \begin{cases}
    0 \le z \le 1+x+y-3(x-y)y \\
    0 \le y \le 1 \\
    y \le x \le y+1
  \end{cases} \label{eq:1}
\end{align}
において，$y=k\,(0\le l \le 1)$で切断した$xz$平面内で考える．
するとこの平面内では\cref{eq:1}は
\begin{align}
  \begin{cases}
    0\le z\le (1-3k)x+(1+k+3k^2)\equiv f(x) \\
    k\le x\le k+1
  \end{cases}
\end{align}
となる．$f(k)=1+2k>0$，$f(k+1)=2-k>0$ゆえ，断面の概形は以下のような台形になる．

\begin{figure}[htb]
  \centering
  \begin{tikzpicture}[scale=2.2, >=stealth, every node/.style={font=\small}]

    % --- 1. 定数・座標パラメータの定義 ---
    % スケッチの比率に合わせた代表値 (例: k = 0.2 付近の見やすい比率)
    % x 範囲: [k, k+1]
    % z 範囲: 左端 2k+1, 右端 2-k
    \def\kVal{1.3}
    \def\wVal{1.7}
    \def\zLeft{1.1}
    \def\zRight{2.2}

    \coordinate (O)     at (0, 0);
    \coordinate (Xk)    at (\kVal, 0);
    \coordinate (Xkone) at ({\kVal + \wVal}, 0);
    \coordinate (Pk)    at (\kVal, \zLeft);
    \coordinate (Pkone) at ({\kVal + \wVal}, \zRight);

    % --- 2. 断面台形領域の斜線ハッチング ---
    \fill[pattern=north east lines, pattern color=blue!65, opacity=0.75]
    (Xk) -- (Pk) -- (Pkone) -- (Xkone) -- cycle;

    % --- 3. 座標軸 (xz 平面) ---
    \draw[->, thick] (-0.3, 0) -- ({\kVal + \wVal + 0.8}, 0) node[right] {$x$};
    \draw[->, thick] (0, -0.3) -- (0, {\zRight + 0.6}) node[above] {$z$};
    \node[below left=2pt] at (O) {$O$};

    % --- 4. 台形の輪郭線 ---
    % 上辺: 直線 z = f(x)
    \draw[very thick, blue!85!black] (Pk) -- (Pkone);
    % 左右の辺
    \draw[thick, black] (Xk) -- (Pk);
    \draw[thick, black] (Xkone) -- (Pkone);
    % 底辺
    \draw[thick, black] (Xk) -- (Xkone);

    % 直線 z = f(x) の延長（参考表示）
    \draw[dashed, blue!50, domain=0.8:3.4]
    plot (\x, {\zLeft + (\zRight - \zLeft)/(\wVal)*(\x - \kVal)});
    \node[blue!85!black, font=\footnotesize, above right] at (Pkone) {$z = f(x)$};

    % --- 5. 座標補助線 (破線) ---
    % z 軸への投影線
    \draw[dashed, gray!70] (Pk) -- (0, \zLeft);
    \draw[dashed, gray!70] (Pkone) -- (0, \zRight);

    % --- 6. 座標目盛りとラベル ---
    % x 軸上の点
    \fill (Xk) circle (0.5pt) node[below=2pt] {$k$};
    \fill (Xkone) circle (0.5pt) node[below=2pt] {$k+1$};

    % z 軸上の目盛り
    \draw[gray!80!black] (-0.05, \zLeft) -- (0.05, \zLeft);
    \node[left=2pt] at (0, \zLeft) {$2k+1$};

    \draw[gray!80!black] (-0.05, \zRight) -- (0.05, \zRight);
    \node[left=2pt] at (0, \zRight) {$2-k$};

    % 頂点のプロット
    \fill (Pk) circle (0.6pt);
    \fill (Pkone) circle (0.6pt);

    % 断面積 S(k) の注釈
    \node[draw=blue!85!black, fill=white, rounded corners=2pt, inner sep=2.5pt, font=\footnotesize]
    at ({\kVal + 0.5*\wVal}, {0.45*(\zLeft + \zRight)}) {$S(k)$};

  \end{tikzpicture}
  \caption{平面 $y = k$ による立体の切断面（台形領域）}
  \label{fig:cross_section_yk}
\end{figure}

したがって，この平面内での題意の立体の断面積$S(k)$は
\begin{align}
  S(y)=\dfrac{1}{2}\left\{(2k+1)+(2-k)\right\}=\dfrac{k+3}{2}
\end{align}
だから，求める体積$V$は
\begin{align}
  V=\int_0^1S(k)dk=\dfrac{1}{2}\left[\dfrac{1}{2}k^2+3k\right]_0^1=\dfrac{7}{4}\tag{答}
\end{align}
である．

\newpage
\end{document}
